% from ega3-1-fr.tex
\oldpage[0\textsubscript{III}]{60}
où les flèches horizontales proviennent de
(\hyperref[0.12.2.2.1-fr]{12.2.2.1}) (la dernière combinée avec
(\hyperref[0.4.2.2.1-fr]{0\textsubscript{I}, 4.2.2.1})) et les flèches
verticales des homomorphismes (\hyperref[0.12.2.5.1-fr]{12.2.5.1}) et
(\hyperref[0.12.2.5.2-fr]{12.2.5.2}) respectivement.

Il suffit en effet de le vérifier pour les homomorphismes correspondants de
préfaisceaux; si on revient aux définitions données dans
(\hyperref[0.12.2.2-fr]{12.2.2}) et
(\hyperref[0.12.2.5-fr]{12.2.5}) pour ces homomorphismes, on est aussitôt
ramené, pour (\hyperref[0.12.2.6.1-fr]{12.2.6.1}), aux diagrammes commutatifs
(\hyperref[0.12.1.5.2-fr]{12.1.5.2}); la vérification est encore plus simple
pour (\hyperref[0.12.2.6.2-fr]{12.2.6.2}).
\end{env}

\subsection{Compléments sur les foncteurs Ext de faisceaux}
\label{subsection:0.12.3-fr}

\begin{env}[12.3.1]
\label{0.12.3.1-fr}
Considérons un espace annelé $(X,\mathcal O_X)$; nous ne reviendrons pas sur
la définition et les principales propriétés des bifoncteurs
$\operatorname{Ext}_{\mathcal O_X}^p(X;\mathcal F,\mathcal G)$ de la
catégorie des $\mathcal O_X$-Modules dans celle des
$\Gamma(X,\mathcal O_X)$-modules, et
$\mathcal{E}xt_{\mathcal O_X}^p(\mathcal F,\mathcal G)$ de la catégorie des
$\mathcal O_X$-Modules dans elle-même, ni sur la suite spectrale birégulière
$E(\mathcal F,\mathcal G)$ qui les relie (T, 4.2 et G, II, 7.3).
\end{env}

\begin{env}[12.3.2]
\label{0.12.3.2-fr}
On définit de la même manière que dans (M, XIV, 1) la notion d'extension d'un
$\mathcal O_X$-Module $\mathcal F$ par un $\mathcal O_X$-Module $\mathcal G$
et la loi de composition entre classes d'extensions équivalentes : les
raisonnements faits pour les modules s'adaptent en effet de façon évidente à
une catégorie abélienne quelconque. La seconde démonstration de (M, XIV, 1.1),
qui n'utilise que l'existence de plongements dans les objets injectifs, est
alors encore valable pour la catégorie des $\mathcal O_X$-Modules, et montre
donc que $\operatorname{Ext}_{\mathcal O_X}^1(X;\mathcal F,\mathcal G)$
s'identifie canoniquement au groupe abélien des classes d'extensions de
$\mathcal F$ par $\mathcal G$.
\end{env}

\begin{proposition}[12.3.3]
\label{0.12.3.3-fr}
\emph{Soit $(X,\mathcal O_X)$ un espace annelé tel que le faisceau d'anneaux
$\mathcal O_X$ soit cohérent. Alors, pour tout couple de
$\mathcal O_X$-Modules cohérents $\mathcal F$, $\mathcal G$ et tout
$p\geq0$, $\mathcal{E}xt_{\mathcal O_X}^p(\mathcal F,\mathcal G)$ est un
$\mathcal O_X$-Module cohérent.}
\end{proposition}

\begin{proof}
Notons que les $\mathcal{E}xt_{\mathcal O_X}^p(\mathcal F,\mathcal G)$
forment un foncteur cohomologique contravariant en $\mathcal F$. Puisque
$\mathcal F$ est cohérent, il existe pour tout $p$ et tout point $x\in X$ un
voisinage ouvert $U$ de $X$ et une suite exacte de
$(\mathcal O_X|U)$-Modules
\[
0\longrightarrow\mathcal R\longrightarrow\mathcal L_{p-1}
\longrightarrow\cdots\longrightarrow\mathcal L_0
\longrightarrow\mathcal F|U\longrightarrow0
\]
où chacun des $\mathcal L_i$ ($0\leq i\leq p-1$) est isomorphe à un
$\mathcal O_X^{n_i}|U$ et $\mathcal R$ est cohérent : cela résulte par
récurrence sur $p$ de (\hyperref[0.5.3.2-fr]{0\textsubscript{I}, 5.3.2}) et
(\hyperref[0.5.3.4-fr]{0\textsubscript{I}, 5.3.4}), vu l'hypothèse que
$\mathcal O_X$ est cohérent.

Notons maintenant que, pour $p\geq1$, on a
\[
\operatorname{Ext}_{\mathcal O_X|U}^p(\mathcal L|U,\mathcal G|U)=0
\]
pour tout $\mathcal O_X$-Module $\mathcal L$ tel que $\mathcal L|U$ soit
isomorphe à un $\mathcal O_X^n|U$ (T, 4.2.3); le raisonnement de (M, V, 7.2)
s'applique donc au foncteur cohomologique contravariant
\[
\mathcal F\leadsto
\mathcal{H}om_{\mathcal O_X|U}(\mathcal F|U,\mathcal G|U),
\]
et donne une suite exacte
\[
\mathcal{H}om_{\mathcal O_X|U}(\mathcal L_{p-1},\mathcal G|U)
\longrightarrow
\mathcal{H}om_{\mathcal O_X|U}(\mathcal R,\mathcal G|U)
\longrightarrow
\mathcal{E}xt_{\mathcal O_X|U}^p(\mathcal F|U,\mathcal G|U)
\longrightarrow0
\]
et comme les deux premiers termes de cette suite sont des
$(\mathcal O_X|U)$-Modules cohérents
(\hyperref[0.5.3.5-fr]{0\textsubscript{I}, 5.3.5}), il en est de même du
troisième (\hyperref[0.5.3.4-fr]{0\textsubscript{I}, 5.3.4}).
\end{proof}

